\chapter{The repair: Theorem $7'$}\label{chap:repair}

This chapter documents the Lean~4 formalization of Theorem~$7'$: the
unconditional repair of the classification of rational-derived quartics
with a repeated root (Theorem~7 of Buchholz--MacDougall,
J.~Number Theory 81 (2000)).  Every node below is formalized; the master
statement \texttt{Thm7Statement.thm7prime} has axiom footprint exactly
\texttt{[propext, Classical.choice, Quot.sound]}.

\section{The normalized family and the gates}

\begin{definition}[Normalized $p(2,1,1)$ family]\label{def:family}
\lean{Thm7Statement.quartic211}\leanok
For $a \in \Q$ let $q_a(x) = x^2(x-1)(x-a)$, the normalized quartic with
root multiset $\{0,0,1,a\}$; it has genuine type $p(2,1,1)$ iff
$a \notin \{0,1\}$.
\end{definition}

\begin{definition}[Gates]\label{def:gates}
\lean{Thm7Statement.Gate1}\leanok
\uses{def:family}
$G_1(a)$: $9a^2-14a+9$ is a rational square.
$G_2(a)$: $9a^2-6a+9$ is a rational square.
(Formalized as \texttt{Gate1} and \texttt{Gate2}.)  These are the
discriminant conditions for $q_a'$ and $q_a''$ to split over $\Q$.
\end{definition}

\begin{definition}[Rational-derivedness, gate form]\label{def:rd}
\lean{Thm7Statement.RD211}\leanok
\uses{def:gates}
$\mathrm{RD}(a)$: $a \neq 0$, $a \neq 1$, $G_1(a)$ and $G_2(a)$.
\end{definition}

\begin{lemma}[Splitting under the gates]\label{lem:splits}
\lean{Thm7Statement.dq1_splits}\leanok
\uses{def:gates}
Under $G_1$ resp.\ $G_2$, the derivatives $q_a'$ and $q_a''$ split with
explicit rational roots.  (Formalized as \texttt{dq1\_splits} and
\texttt{dq2\_splits}.)
\end{lemma}
\begin{proof}\leanok
Quadratic formula; a one-line \texttt{linear\_combination}.
\end{proof}

\begin{lemma}[Normalization of a general $p(2,1,1)$ translate]
\label{lem:normalize}
\lean{Thm7Statement.caseA_normalized}\leanok
\uses{def:family}
For $r_1 \neq c$, the quartic $(x-c)^2(x-r_1)(x-r_2)$ pulled back along
$x = c + (r_1-c)u$ equals $(r_1-c)^4\, q_a(u)$ with
$a = (r_2-c)/(r_1-c)$.
\end{lemma}
\begin{proof}\leanok
A ring identity.
\end{proof}

\section{The curve and the $a$-map}

\begin{definition}[The curve $E$]\label{def:E}
\lean{Thm7Statement.OnE}\leanok
$E\colon z^2 = w^3+12w^2-108w = w(w-6)(w+18)$, Cremona \texttt{576i2}.
The membership predicate is \texttt{OnE}; the mathlib Weierstrass model
is \texttt{E576i2}, tied to \texttt{OnE} by
\texttt{OnE\_iff\_equation}, with $\Delta \neq 0$ proved in
\texttt{E576i2\_elliptic}.
\end{definition}

\begin{definition}[The $a$-map]\label{def:amap}
\lean{Thm7Statement.aMap}\leanok
\uses{def:E}
$a(w,z) = \dfrac{9(2w+z-12)(w+2)}{(z-w-18)(8w+z)}$, written with the
on-curve reduced denominator
$D(w,z) = (7w-18)z + w^3+4w^2-252w$ (\texttt{D\_reduce},
\texttt{aNum}, \texttt{aDen}).
\end{definition}

\section{The statement}

\begin{definition}[Theorem $7'$, the statement]\label{def:thm7prime}
\lean{Thm7Statement.Thm7Prime}\leanok
\uses{def:rd, def:amap}
$\mathrm{Thm7'}$ is the conjunction of the backward direction
(\texttt{Thm7Backward}: every $a$-map value on $E(\Q)$ away from the
denominator locus and from $\{0,1\}$ satisfies $\mathrm{RD}$) and the
forward direction (\texttt{Thm7Forward}: every $a$ with $\mathrm{RD}(a)$
arises as such a value).
\end{definition}

\section{The backward direction}

\begin{lemma}[Cleared gate identities]\label{lem:gatescleared}
\lean{Thm7Statement.gate1_cleared}\leanok
\uses{def:E, def:amap}
On $E$, cleared of denominators, $9a^2-14a+9$ and $9a^2-6a+9$
(for $a = a(w,z)$) are squares of explicit rational functions;
the certificates are single \texttt{linear\_combination} steps
(\texttt{gate1\_cleared}, \texttt{gate2\_cleared}), packaged in division
form as \texttt{gate1\_sq}, \texttt{gate2\_sq}.
\end{lemma}
\begin{proof}\leanok
Sage-lifted cofactor certificates, checked by the kernel.
\end{proof}

\begin{lemma}[Exceptional locus]\label{lem:exceptional}
\lean{Thm7Statement.exceptional_points}\leanok
\uses{def:E, def:amap}
On $E(\Q)$, the side factor $w(w+18)(w^2-63w+486)$ vanishes with
$D(w,z) \neq 0$ only at $(54,432)$ and $(9,-27)$; both map to
$a = 77/90$, where the gates hold with witnesses $19/10$, $97/30$
(\texttt{gates\_at\_77\_90}).
\end{lemma}
\begin{proof}\leanok
Finite case analysis.
\end{proof}

\begin{theorem}[Backward direction]\label{thm:backward}
\lean{Thm7Statement.thm7_backward}\leanok
\uses{def:thm7prime, lem:gatescleared, lem:exceptional}
\texttt{Thm7Backward} holds.
\end{theorem}
\begin{proof}\leanok
Generic branch by \Cref{lem:gatescleared}, exceptional branch by
\Cref{lem:exceptional}.  No rank input.
\end{proof}

\section{The forward direction}

\begin{definition}[The inverse map]\label{def:inverse}
\lean{Thm7Statement.wNp}\leanok
\uses{def:gates, def:E}
Given gate witnesses $(r,s)$ for $a$, set
$w = 18(15a-6r-s-13)/(3r-7s+12)$ and determine $z$ linearly from
$a \cdot D(w,z) = 9(2w+z-12)(w+2)$; the denominators involve
$K(a,r,s) = 105a^2-45ar-238a+51r+16s+105$ and the sign-selection sextic
$P(a,r,s)$ (\texttt{wNp}, \texttt{wDp}, \texttt{zNp}, \texttt{zDp},
\texttt{K1p}, \texttt{adP6}).  Discovered by Gr\"obner elimination.
\end{definition}

\begin{lemma}[Cleared inverse identities]\label{lem:inversecleared}
\lean{Thm7Statement.curve_cleared}\leanok
\uses{def:inverse}
Modulo the two gate relations, the inverse point lies on $E$
(\texttt{curve\_cleared}) and the pulled-back denominator equals
$-62208\,P(a,r,s)$ up to the chart clearing (\texttt{aden\_cleared}).
\end{lemma}
\begin{proof}\leanok
Sage-lifted cofactor certificates, checked by the kernel.
\end{proof}

\begin{lemma}[$96$ is not a rational square]\label{lem:sq96}
\lean{Thm7Statement.sq_ne_96}\leanok
For all $x \in \Q$, $x^2 \neq 96$.
\end{lemma}
\begin{proof}\leanok
Via \texttt{Rat.sqrt}.
\end{proof}

\begin{lemma}[Denominator nonvanishing, all signs]\label{lem:denoms}
\lean{Thm7Statement.wD_ne}\leanok
\uses{def:inverse, lem:sq96}
If the gates hold and $a \neq 0$, then $3r-7s+12 \neq 0$ and
$K(a,r,s) \neq 0$ for every choice of signs of $(r,s)$: the eliminants
are $a^2(225a^2-210a+238)$ with $225a^2-210a+238 = (15a-7)^2+189 > 0$,
and $(3a-5)(5a-3)a^2(225a^2-210a+238)$, whose linear branches force
$(3r)^2 = 96$ resp.\ $(5r)^2 = 96$
(\texttt{wD\_ne}, \texttt{K1\_ne}).
\end{lemma}
\begin{proof}\leanok
Eliminant certificates plus \Cref{lem:sq96} and positivity.
\end{proof}

\begin{lemma}[Sign cover]\label{lem:cover}
\lean{Thm7Statement.adP6_cover}\leanok
\uses{def:inverse}
Under the gates, at least one of the four sign choices keeps
$P(a,\pm r,\pm s) \neq 0$: the ideal generated by the gate relations and
all four conjugates is the unit ideal, and the Nullstellensatz
certificate is checked by the kernel.
\end{lemma}
\begin{proof}\leanok
A single \texttt{linear\_combination} deriving $1 = 0$ from the contrary.
\end{proof}

\begin{lemma}[Forward core]\label{lem:core}
\lean{Thm7Statement.forward_core}\leanok
\uses{lem:inversecleared, lem:denoms}
Under the gates and the three nonvanishing conditions, the inverse point
lies on $E(\Q)$, avoids the $a$-map denominator locus, and maps back to
$a$.
\end{lemma}
\begin{proof}\leanok
Assemble \Cref{lem:inversecleared} with the chart clearings.
\end{proof}

\begin{theorem}[Forward direction]\label{thm:forward}
\lean{Thm7Statement.thm7_forward}\leanok
\uses{def:thm7prime, lem:core, lem:cover, lem:denoms}
\texttt{Thm7Forward} holds.  No rank input.
\end{theorem}
\begin{proof}\leanok
Choose signs by \Cref{lem:cover}; apply \Cref{lem:core}.
\end{proof}

\section{The main theorem}

\begin{theorem}[Theorem $7'$, unconditional]\label{thm:main}
\lean{Thm7Statement.thm7prime}\leanok
\uses{thm:backward, thm:forward}
$\mathrm{Thm7'}$ holds, with axiom footprint exactly
\texttt{[propext, Classical.choice, Quot.sound]}.
\end{theorem}
\begin{proof}\leanok
$\langle$\texttt{thm7\_backward}, \texttt{thm7\_forward}$\rangle$.
\end{proof}

\begin{remark}[The disclosed rank axiom]\label{rem:rank}
\uses{thm:main}
The development retains one disclosed, instance-localized axiom
(\texttt{rank\_E576i2}: the Mordell--Weil rank of \texttt{576i2} is $1$,
verified unconditionally in two independent systems), consumed only by
two auxiliary declarations from an earlier conditional architecture.
The main theorem does not depend on it, and the kernel's
per-declaration axiom report documents this separation.
\end{remark}

\section{Statement fidelity in the kernel}

The gate form (Definition~\ref{def:rd}) was chosen at design time as the
division-free, verification-friendly restatement of rational-derivedness; its
correspondence with the natural form was originally carried by the CAS
statement-fidelity gate (\texttt{scripts/sage/wp1\_fidelity.sage}).  The two
links are now closed in the kernel as well
(\texttt{Thm7Prime/Fidelity.lean}); the CAS gate remains as the independent
design-time cross-check.

\begin{definition}[Natural form]\label{def:naturalrd}\lean{Thm7Fidelity.NaturalRD}\leanok
\uses{def:rd}
$\mathrm{NaturalRD}(a)$: the polynomial $Q_a = X^{2}(X-1)(X-\mathrm{C}\,a)$ and
its first three \texttt{Polynomial.derivative}s split over $\Q$ (the fourth
derivative is constant), mirroring the transcript's \texttt{KDerived}.
\end{definition}

\begin{theorem}[Gate form $\iff$ natural form]\label{thm:fidelity}\lean{Thm7Fidelity.rd211_iff_natural}\leanok
\uses{def:naturalrd, def:rd}
$\mathrm{RD211}(a) \iff a\neq 0 \wedge a\neq 1 \wedge \mathrm{NaturalRD}(a)$.
\end{theorem}
\begin{proof}\leanok
Forward: the kernel-computed derivatives (\texttt{derivative\_Q},
\texttt{derivative2\_Q}, with evaluation anchors to \texttt{dq1},
\texttt{dq2}) split explicitly under the gates
(\texttt{derivative\_splits\_of\_gate1},
\texttt{derivative2\_splits\_of\_gate2}); the quartic and the linear third
derivative split outright.  Converse: a quadratic over $\Q$ with nonzero
leading coefficient that splits has square discriminant
(\texttt{disc\_sq\_of\_quadratic\_splits}), applied to both derivatives.
\end{proof}

\section{The affine classification}

The reviewer-facing form of the result: the bridge from an arbitrary split
$(2,1,1)$ quartic to the normalised family, packaged as one kernel theorem
(\texttt{Thm7Prime/Classification.lean}).

\begin{definition}[The equivalence induced by $\langle X^{*}\rangle$]\label{def:affequiv}\lean{Thm7Classification.AffineEquiv}\leanok
$f\sim g$ iff $g=\nu\,f(\lambda X+\mu)$ with $\lambda,\nu\neq 0$.  This is an
equivalence relation (\texttt{affineEquiv\_refl}, \texttt{affineEquiv\_symm},
\texttt{affineEquiv\_trans}, all kernel-checked).
\end{definition}

\begin{definition}[Split $(2,1,1)$ quartic]\label{def:isp211}\lean{Thm7Classification.IsP211}\leanok
$f=\lambda(X-c)^{2}(X-r_{1})(X-r_{2})$ with $\lambda\neq0$ and $c,r_{1},r_{2}$
pairwise distinct.
\end{definition}

\begin{definition}[Rational-derivedness, natural form]\label{def:rdpoly}\lean{Thm7Classification.RDPoly}\leanok
\uses{def:naturalrd}
$f$ and its first three derivatives split over $\Q$ (for quartics, exactly
rational-derivedness; used only under Definition~\ref{def:isp211}).
\end{definition}

\begin{theorem}[Affine classification]\label{thm:classification}\lean{Thm7Classification.classification}\leanok
\uses{def:affequiv, def:isp211, def:rdpoly, thm:fidelity}
For a split $(2,1,1)$ quartic $f$: $f$ is rational-derived iff $f$ is
$\langle X^{*}\rangle$-equivalent to $Q_{a}$ for some $a$ with
$\mathrm{RD211}(a)$.
\end{theorem}
\begin{proof}\leanok
Rational-derivedness is $\langle X^{*}\rangle$-invariant (chain rule and
transport of splitting along linear substitutions); the normalisation
$a=(r_{2}-c)/(r_{1}-c)$ is the kernel form of \texttt{caseA\_normalized};
the gates then come from Theorem~\ref{thm:fidelity}.
\end{proof}

\begin{theorem}[Classification by the curve]\label{thm:classcurve}\lean{Thm7Classification.classification_by_curve}\leanok
\uses{thm:classification, thm:main}
For a split $(2,1,1)$ quartic $f$: $f$ is rational-derived iff $f$ is
$\langle X^{*}\rangle$-equivalent to $Q_{a\mathrm{Map}(w,z)}$ for an affine
rational point $(w,z)$ of the Weierstrass model of $E=\texttt{576i2}$ with
$a\mathrm{Den}(w,z)\neq0$ and $a\mathrm{Map}(w,z)\notin\{0,1\}$.
\end{theorem}
\begin{proof}\leanok
Theorem~\ref{thm:classification} composed with both directions of the main
theorem (\texttt{thm7prime}).
\end{proof}
